\documentclass[11pt,a4paper]{article}
\usepackage[margin=2.5cm]{geometry}
\usepackage[T1]{fontenc}
\usepackage{booktabs,longtable,array}
\usepackage{listings}
\usepackage{xcolor}
\usepackage{hyperref}
\usepackage{enumitem}

\definecolor{codebg}{gray}{0.95}
\lstset{basicstyle=\ttfamily\footnotesize,backgroundcolor=\color{codebg},breaklines=true,
        breakatwhitespace=false,frame=single,columns=fullflexible,keepspaces=true,showstringspaces=false}

\newcommand{\repourl}{\url{https://github.com/rubimburagad85-hash/cryptography-network-security-exam}}

\title{Cryptography and Network Security\\ \large Integrated Situation --- Technical Report\\
        \normalsize ETTCS801 \,|\, ULK Polytechnic Institute}
\author{Gad RUBIMBURA \\ Registration Number: 4202670049}
\date{\today}

\begin{document}
\maketitle
\tableofcontents
\bigskip

%---------------------------------------------------------------------
\section{Introduction}
A polytechnic institute stores student records on a central server and transfers files between two campuses[cite: 1].
A security review found weak staff passwords, outdated software, unencrypted file transfers and guest-network
access to the records server, and repeated connection attempts from an unfamiliar external address were observed[cite: 1].
This report documents the risk assessment, the Python encryption and integrity toolkit, and the firewall
rules produced to address these findings[cite: 1]. All source files are in the GitHub repository (Section~\ref{sec:conclusion}).

%---------------------------------------------------------------------
\section{Risk Assessment}
\subsection{Assets, vulnerabilities and consequences}
\begin{longtable}{@{}p{0.3cm}p{3.3cm}p{4.2cm}p{6.3cm}@{}}
\toprule
\# & Asset & Vulnerability & Possible consequence\\ \midrule\endhead
R1 & Staff accounts and credentials & Weak passwords; repeated external login attempts &
Account takeover; records read, altered or deleted; misuse under a staff identity[cite: 1].\\[4pt]
R2 & Student records server & Guest network can reach the server; outdated software &
Scanning and exploitation of known flaws; data breach, ransomware, outage; reputational and legal harm[cite: 1].\\[4pt]
R3 & Record files in transit between campuses & Unencrypted file transfer &
Eavesdropping or man-in-the-middle; records disclosed or silently modified[cite: 1].\\
\bottomrule
\end{longtable}

\subsection{Risk ranking}
Likelihood and impact are scored from 1 (low) to 5 (high); the score is their product[cite: 1].
\begin{center}
\begin{tabular}{@{}clccc@{}}
\toprule
Rank & Risk & Likelihood & Impact & Score\\ \midrule
1 & R1 -- Weak staff passwords & 5 & 5 & 25\\
2 & R2 -- Guest access / outdated software & 4 & 5 & 20\\
3 & R3 -- Unencrypted transfers & 3 & 4 & 12\\ \bottomrule
\end{tabular}
\end{center}
\textbf{Reasoning.} R1 ranks first because attacks are already occurring and weak passwords need little skill to
defeat; one compromised account exposes all data the user can reach[cite: 1]. R2 is second: guests are many and untrusted, and
known exploits exist for outdated software, with very high impact because all records reside on one server[cite: 1]. R3 ranks
third because interception requires a position on the inter-campus path, although success exposes and allows tampering with complete files[cite: 1].

\subsection{Recommended controls}
\begin{center}
\begin{tabular}{@{}lp{11cm}@{}}
\toprule
Risk & Control\\ \midrule
R1 & Strong password policy (minimum 12 characters, no reuse), multi-factor authentication, account lockout / rate limiting[cite: 1].\\
R2 & Network segmentation with firewall rules denying guest access to the server, plus a regular patching schedule[cite: 1].\\
R3 & Encrypt files before transfer (AES-256-GCM) and use secure channels (SFTP/TLS/VPN), with SHA-256 integrity checks[cite: 1].\\
\bottomrule
\end{tabular}
\end{center}

%---------------------------------------------------------------------
\section{Encryption and Integrity Application}
\subsection{Design}
The programme \texttt{src/secure\_records.py} is written in Python and uses the \texttt{cryptography} library[cite: 1].
\begin{itemize}[nosep]
  \item \textbf{Encryption} -- AES-256-GCM authenticated encryption. A random 256-bit key is generated once and stored
        \emph{outside} the repository (default \texttt{\textasciitilde/.ulk\_keys/records.key}); the tool refuses a key path inside a git
        repository[cite: 1]. Each encryption uses a fresh random 96-bit nonce[cite: 1]. The output file is
        \texttt{header (4 bytes) \(\|\) nonce (12 bytes) \(\|\) ciphertext \(\|\) GCM tag (16 bytes)}[cite: 1].
  \item \textbf{Decryption} -- reads the nonce, decrypts and authenticates[cite: 1]. A wrong key or any modified byte raises an
        authentication failure, reported as a clean error message[cite: 1]. With \texttt{--original} the decrypted file's SHA-256
        is compared with the original's to prove the contents match[cite: 1].
  \item \textbf{Integrity} -- \texttt{hash} stores the SHA-256 digest of a file in a baseline file (\texttt{hashes.json});
        \texttt{check} recomputes the digest and reports \emph{UNCHANGED} or \emph{MODIFIED}[cite: 1].
  \item \textbf{Error handling} -- missing files, directories, empty files, invalid or missing keys, corrupt or non-encrypted
        input and corrupt baselines are caught and reported with a message and exit code, never a traceback[cite: 1].
\end{itemize}

\subsection{Test evidence}
Ten automated unit tests (round trip, ciphertext differs from plaintext, tamper detection, wrong key, missing key,
missing/empty/invalid input, hash change detection, CLI error code) all pass[cite: 1]:
\begin{lstlisting}
$ python3 -m unittest discover -s tests -v
...
Ran 10 tests in 0.008s
OK
\end{lstlisting}
Full demonstration output (\texttt{tests/run\_demo.sh}); the two SHA-256 values match after decryption, and the
modified file is flagged[cite: 1]:
{\footnotesize
\lstinputlisting[basicstyle=\ttfamily\tiny]{../tests/crypto_demo_output.txt}}

%---------------------------------------------------------------------
\section{Network Traffic Filtering}
\subsection{Laboratory setup}
The laboratory environment uses the records server IP \texttt{192.168.100.10}, authorised staff network \texttt{192.168.10.0/24}, guest network \texttt{10.0.99.0/24}, firewall tool \texttt{iptables}, and the tested service port is \texttt{tcp/443}[cite: 1].

\subsection{Firewall rules}
Rules are in \texttt{firewall/records\_server\_firewall.sh} and use a dedicated iptables chain so they can be removed cleanly[cite: 1].
In order of evaluation:
\begin{enumerate}[nosep]
  \item Accept established/related traffic (and loopback) so existing sessions keep working[cite: 1].
  \item \textbf{Block guest \(\rightarrow\) server (all ports):} \texttt{-s GUEST\_NET -j DROP} (logged, rate limited)[cite: 1].
  \item \textbf{Permit staff \(\rightarrow\) named service:} \texttt{-s STAFF\_NET -p tcp --dport PORT -j ACCEPT}[cite: 1].
  \item \textbf{Block all other inbound access to the service:} \texttt{-p tcp --dport PORT -j DROP}[cite: 1].
\end{enumerate}
\begin{lstlisting}
iptables -A RECORDS_FILTER -m conntrack --ctstate ESTABLISHED,RELATED -j ACCEPT
iptables -A RECORDS_FILTER -s $GUEST_NET -j DROP
iptables -A RECORDS_FILTER -s \(STAFF_NET -p tcp --dport\)SERVICE_PORT \
         -m conntrack --ctstate NEW -j ACCEPT
iptables -A RECORDS_FILTER -p tcp --dport $SERVICE_PORT -j DROP
iptables -A RECORDS_FILTER -j RETURN
iptables -I INPUT 1 -j RECORDS_FILTER
\end{lstlisting}

\subsection{Test procedure and results}
One permitted and two blocked connections were tested with \texttt{nc -zv -w 5 SERVER\_IP PORT}[cite: 1].
\begin{center}
\begin{tabular}{@{}llp{3.3cm}p{4.5cm}@{}}
\toprule
Test & Source & Expected & Actual\\ \midrule
T1 permitted & Staff client & Connection succeeds & Connection to 192.168.100.10 port 443 succeeded!\\
T2 blocked & Guest client & Timeout / dropped & nc: connect to 192.168.100.10 port 443 timed out\\
T3 blocked & Other client & Timeout / dropped & nc: connect to 192.168.100.10 port 443 timed out\\ \bottomrule
\end{tabular}
\end{center}

%---------------------------------------------------------------------
\section{Conclusion}\label{sec:conclusion}
The assessment ranked weak passwords, guest access to the records server, and unencrypted transfers as the main risks[cite: 1].
The toolkit addresses the last two directly: authenticated AES-256-GCM encryption protects files in transit and at rest, SHA-256
baselines detect tampering, and the firewall removes the guest path to the server while limiting the service to the staff
network[cite: 1]. Weak passwords are addressed by policy and MFA, which are organisational controls outside the code[cite: 1]. Limitations: the key must
be distributed securely between campuses (a key-management system would be needed at scale), and firewall rules cover only
the tested service[cite: 1]. Future work: automated patching, centralised logging and intrusion detection[cite: 1].

\medskip
\noindent\textbf{GitHub repository:} \repourl

\end{document}
